% LaTeX-ready appendix fragment, not a standalone paper.
% Requires amsmath, amssymb, amsthm, hyperref in the parent document.
% The results are scoped analyses, not seven claimed novel theorems.
\providecommand{\acamE}{\mathbb{E}}
\providecommand{\acamP}{\mathbb{P}}
\newtheorem{acamprop}{Proposition}

\section{Scoped analysis of imperfect concept acquisition}
\label{app:acam-theory}

We give seven finite-scope results that specify implementation invariants,
analytical test cases, and statistical guardrails. We do not claim that
these constitute six novel theorems. In particular, the noiseless XOR
obstruction and coalition lookahead already appear in RouteCert,
Proposition S2; finite-family risk validation is a standard
learn-then-test construction. Our analysis parameterizes measurement
noise and non-additive execution costs to obtain falsifiable experimental
predictions. It does not establish superiority on real datasets.

\paragraph{Setup.}
The measurement-item/query-group bank is fixed and finite, with group
indices $g\in[K]$ and atomic-attribute indices $j\in[D]$; in general
$K\neq D$. Each group response $Z_g$ has a finite typed space, which may
be a finite product space of atomic responses (a tuple or vector).
Acquisition counts and costs below use these measurement items/groups;
the XOR example is the special case of one bit per group.
For a data unit $W=(X,C,Y)\sim P$, a responder alone may
read $X$. It returns automatic measurements, not assumed to equal $C$.
After $t$ calls the observed transcript is
$H_t=((A_1,Z_{A_1}),\ldots,(A_t,Z_{A_t}))$, with
$H_0=\varnothing$. The head reads the canonical observed set $\bar H_t$.
The controller and candidate generator read only the transcript and
fixed public parameters. A missing annotation is not a runtime response.
When batching changes a responder's output, the complete batch, prompt,
and decoding configuration are part of the action identity.

The action set includes \textnormal{\textsc{Stop}}, which terminates acquisition and
returns a classification; it does not mean abstention. Bounded-loss
results below require $0\leq\ell\leq1$. Ordinary unbounded cross-entropy
does not satisfy this assumption. All results are non-asymptotic.

\begin{acamprop}[T1: structural transcript noninterference]
\label{prop:acam-noninterference}
Fix all learned parameters and public configuration. Suppose
\[
 A_{t+1}=\pi(H_t,b(H_t),U_{t+1}),\qquad
 \widehat Y_t=f(\bar H_t),
\]
where $U_{t+1}$ is case-independent exogenous randomness and the budget
$b(H_t)$ is computed from the observed transcript and a fixed public
cost rule. Neither function nor candidate generation reads $X$,
unqueried responses, labels, input length, measured latency, or
annotation-missingness metadata.

For every legal transcript $h$, every pair $x,x'$, and every common
random seed $u$, replaying $h$ yields identical next actions and current
predictions before the next genuine responder call. If future responder
outputs are also coupled to be identical by a replay hook, the complete
remaining action, stopping, and prediction trajectories are identical.
Nevertheless, the terminal acquisition mask may perfectly predict $Y$.
\end{acamprop}

\begin{proof}
At the same $h$, the budget, candidate set, and head output are equal
because they are the same functions evaluated at the same arguments.
Using the same $u$ makes the controller output equal. If the action is
\textnormal{\textsc{Stop}}, replay terminates identically. Otherwise a common replay
response produces the same updated transcript. Finite induction proves
the complete-trajectory statement. Integrating over identical
case-independent randomization gives equality of replay distributions.
This is a structural replay statement, not an unqualified observational
conditional-independence claim.

For the mask counterexample, let $C_1=Y\sim\operatorname{Bernoulli}(1/2)$
and measure $C_1$ without error. Always query concept 1 first. If its value
is 1, query an irrelevant concept 2 and stop; otherwise stop immediately.
All routing uses acquired evidence only. Yet the terminal mask is
$\{1,2\}$ exactly when $Y=1$, so its indicator predicts $Y$ perfectly.
\end{proof}

This result does not establish the semantic truth of responses or causal
interpretability. A genuine next measurement may change after replacing
$x$; only the pre-measurement controller and head are held invariant.
If measured latency or token counts update a controller-visible budget,
they form additional input-dependent channels and must either be
excluded from the strict variant or explicitly audited.

\begin{acamprop}[T2: finite-action decision stability]
\label{prop:acam-regret}
At a fixed reachable history $h$, let $\mathcal A(h)$ be finite and
nonempty. For each action $a$, let $r_a\in[0,1]$ be the true expected loss
after taking $a$ and immediately classifying with a fixed head, and let
$\kappa_a\geq0$ be its expected incremental cost. The incremental cost of
\textnormal{\textsc{Stop}} is zero. Fix $\lambda\geq0$ and assume
\[
 \max_a|\widehat r_a-r_a|\leq\epsilon,\qquad
 \max_a|\widehat\kappa_a-\kappa_a|\leq\eta.
\]
If $\widehat a$ minimizes the estimated risk-plus-cost objective up to
$\xi\geq0$, and $a^*$ minimizes the true objective on the same action
set, then
\[
 0\leq
 (r_{\widehat a}+\lambda\kappa_{\widehat a})
 -(r_{a^*}+\lambda\kappa_{a^*})
 \leq2\epsilon+2\lambda\eta+\xi.
\]
\end{acamprop}

\begin{proof}
Write $J_a=r_a+\lambda\kappa_a$,
$\widehat J_a=\widehat r_a+\lambda\widehat\kappa_a$, and
$d=\epsilon+\lambda\eta$. The triangle inequality gives
$|J_a-\widehat J_a|\leq d$ simultaneously for all actions. Consequently,
\[
 J_{\widehat a}\leq\widehat J_{\widehat a}+d
 \leq\widehat J_{a^*}+d+\xi
 \leq J_{a^*}+2d+\xi.
\]
The lower bound follows from the optimality of $a^*$. If the unique
optimal action has gap greater than $2d+\xi$, the same bound rules out
selection of a suboptimal action.
\end{proof}

\paragraph{Conditional finite-horizon extension.}
Suppose there are at most $T$ decisions, with zero-cost absorbing states
after stopping. Let $Q_t^*(h,a)$ be the true Bellman action value for the
same terminal loss plus $\lambda$ times all future costs. A reachable
history here means one reachable from the same initial conditions under
the true transition mechanism by any legal action sequence or policy
in the compared action class. In particular, this includes every
history visited by the learned policy $\widehat\pi$, not merely the
support of training histories or an oracle policy. If
$|\widehat Q_t(h,a)-Q_t^*(h,a)|\leq d_t$ for every such history and every
legal action there, and the policy is $\xi_t$-greedy with respect to
$\widehat Q_t$, then
\[
 J^{\widehat\pi}(H_0)-J^{\pi^*}(H_0)
 \leq\sum_{t=1}^T(2d_t+\xi_t).
\]
Indeed, the preceding argument yields
$Q_t^*(h,\widehat a_t)-V_t^*(h)\leq2d_t+\xi_t$. With
$D_t(h)=V_t^{\widehat\pi}(h)-V_t^*(h)$, Bellman recursion gives
\[
 D_t(h)\leq2d_t+\xi_t+
 \acamE[D_{t+1}(H_{t+1})\mid h,\widehat a_t].
\]
Backward induction from $D_{T+1}=0$ proves the claim. A one-step risk
model is not $Q_t^*$, and low average Brier loss does not establish these
uniform assumptions. No finite-sample learning bound for $d_t$ is claimed.

\begin{acamprop}[T3: execution-cost feasibility]
\label{prop:acam-cost}
Fix a backend, precision, input size, concurrency, and caching protocol,
and use a single common cost unit. Suppose the full-query baseline costs
$T_{\rm full}=s+d_K+h$, where $s$ is shared encoding, $d_K$ is the full
response cost, and $h$ is the final head cost. An adaptive run makes
$\tau$ calls, pays nonnegative response costs $d(A_t)$, the same final
$h$, and nonnegative extra controller, intermediate-head, and
communication cost $O$. Assume finite expectations.

If encoding is cached once whenever $\tau\geq1$, with
$p_0=\acamP(\tau=0)$, then
\[
 T_{\rm full}-\acamE T_{\rm adapt}
 =p_0s+d_K-\acamE\sum_{t=1}^{\tau}d(A_t)-\acamE O.
\]
If every call re-encodes the input, then
\[
 T_{\rm full}-\acamE T_{\rm adapt}
 =(1-\acamE\tau)s+d_K-\acamE\sum_{t=1}^{\tau}d(A_t)-\acamE O.
\]
For a nonempty deterministic run with per-call setup cost $a$, per-concept
cost $d$, per-call controller cost $o$, total queried count $q$, and
$\tau\geq1$ calls, adaptive execution is strictly cheaper than one full
batch if and only if
\[
 d(K-q)>a(\tau-1)+o\tau.
\]
Here common encoding and final-head costs cancel.
\end{acamprop}

\begin{proof}
In the cached case,
$T_{\rm adapt}=\mathbf1\{\tau\geq1\}s+\sum_td(A_t)+O+h$.
Take expectations and subtract from $s+d_K+h$. Replacing the first term
by $\tau s$ proves the re-encoding identity. For the last claim subtract
$s+a\tau+dq+o\tau+h$ from $s+a+dK+h$ and rearrange.
\end{proof}

In particular, if $p_0=0$ and expected response plus extra overhead is at
least $d_K$, no average saving is possible in this model. Even with free
responses and control, the cached speedup is at most
$(s+d_K+h)/(s+h)$ when $s+h>0$. An empty-history stop cannot selectively
identify easy cases without previously paid evidence or an additional
input channel. These are accounting statements within a specified cost
model, not universal computational lower bounds. Asynchronous overlap,
dynamic batching, and cache-dependent answers require fresh measurements.

\begin{acamprop}[T4: imperfect complementary measurements]
\label{prop:acam-xor}
Let $C_1,C_2$ be independent uniform bits, $Y=C_1\oplus C_2$, and
$Z_j=C_j\oplus E_j$, where
$E_1,E_2\stackrel{\rm iid}{\sim}\operatorname{Bernoulli}(\rho)$ are
independent of $(C_1,C_2)$ and $0\leq\rho\leq1/2$. Measurements are fixed
once per case, and classification is Bayes-optimal under zero-one loss.
Then
\[
 R_\varnothing=R_{\{1\}}=R_{\{2\}}=\frac12,\qquad
 R_{\{1,2\}}=2\rho(1-\rho),
\]
and the joint acquisition value is
\[
 \Delta_{12}=\frac12(1-2\rho)^2.
\]
A myopic exact risk-plus-cost controller restricted to single queries
and \textnormal{\textsc{Stop}} stops immediately if both single-query penalties are
strictly positive. A pair action of cost $\kappa_{12}$ is strictly better
than stopping precisely when
\[
 \lambda\kappa_{12}<\frac12(1-2\rho)^2.
\]
Writing $v=\lambda\kappa_{12}$, the favorable region for $0\leq v<1/2$ is
$0\leq\rho<(1-\sqrt{2v})/2$; no strict improvement is possible for
$v\geq1/2$. Equalities describe ties.
\end{acamprop}

\begin{proof}
The unobserved label is uniform. Conditioning on $Z_1$ leaves $C_2$ an
independent uniform bit, so $Y$ remains uniform; the same holds for
$Z_2$. Thus the empty and singleton Bayes errors equal $1/2$.

Let $E=(E_1,E_2)$ and $Z=(Z_1,Z_2)$. For every $e,z\in\{0,1\}^2$,
\[
 \acamP(E=e,Z=z)=\acamP(E=e)\acamP(C=z\oplus e)
 =\tfrac14\acamP(E=e).
\]
Hence $E$ is independent of $Z$. The parity noise
$N=E_1\oplus E_2$ has probability
$q=2\rho(1-\rho)\leq1/2$, and
$Y=Z_1\oplus Z_2\oplus N$. Conditional on $Z$, predicting the observed
parity is Bayes-optimal and errs with probability $q$. Subtracting $q$
from $1/2$ gives the stated value.

Each singleton has objective $1/2+\lambda\kappa_j>1/2$, whereas stopping
has objective $1/2$. The pair has objective
$2\rho(1-\rho)+\lambda\kappa_{12}$; comparison with $1/2$ gives the
criterion. Because $1-2\rho\geq0$, taking the nonnegative square root
gives the threshold. The maximal gain is $1/2$, proving the remaining
boundary cases.
\end{proof}

\paragraph{Correlated-error extension.}
Keep $E$ independent of the uniform concept vector but allow dependence
between $E_1$ and $E_2$. The preceding factorization still holds. With
$d=\acamP(E_1\neq E_2)$, the Bayes pair error is
$\min(d,1-d)$ and the value is $1/2-\min(d,1-d)$. For identical marginal
error $\rho$, let $u=\acamP(E_1=E_2=1)$; then $d=2\rho-2u$.
At $\rho=1/4$, independent errors give pair error $3/8$, perfectly
same-direction errors give zero, and mutually exclusive errors give
$1/2$. All have marginal concept accuracy $3/4$. Thus error dependence
interacts with the task; stronger correlation is not uniformly harmful.

This construction does not defeat every sequential policy: two-step
lookahead, an optimal Bellman policy, and forced acquisition of both
features can succeed. It also does not prove an instance-adaptive
advantage: a static pair attains exactly the same risk as the pair-aware
controller. Those controls are required in the mechanism experiment.

\begin{acamprop}[T5: finite-family terminal-risk validation]
\label{prop:acam-certification}
Condition on training and design data. Fix $M\geq1$ complete systems
before observing an independent calibration sample of $n\geq1$ iid
deployment-distributed cases. Each system includes its responder, head,
controller, stopping thresholds, and failure handling. Randomized
rollouts use independent deployment-distributed randomness across cases;
different systems may be dependent on the same case. Each complete case
produces one terminal loss $L_{i,m}\in[0,1]$. Define
\[
 R_m=\acamE L_{i,m},\quad
 \widehat R_m=\frac1n\sum_{i=1}^nL_{i,m},\quad
 \gamma=\sqrt{\frac{\log(M/\delta)}{2n}},\quad 0<\delta<1,
\]
\[
 U_m=\min\{1,\widehat R_m+\gamma\},\qquad
 \mathcal C_\alpha=\{m:U_m\leq\alpha\},\quad 0\leq\alpha\leq1.
\]
With probability at least $1-\delta$, all $R_m\leq U_m$ simultaneously.
For any data-dependent choice $\widehat m\in\mathcal C_\alpha$ whenever
that set is nonempty,
\[
 \acamP\{\mathcal C_\alpha\neq\varnothing
          \text{ and }R_{\widehat m}>\alpha\}\leq\delta.
\]
\end{acamprop}

\begin{proof}
For completeness, let $X\in[0,1]$ and
$\psi(t)=\log\acamE e^{t(X-\acamE X)}$. Boundedness permits two
derivatives. Under exponential tilting,
$\psi''(t)=\operatorname{Var}_t(X)\leq1/4$, since every distribution
supported on $[0,1]$ satisfies
$\operatorname{Var}(X)\leq\acamE(X-1/2)^2\leq1/4$.
Together with $\psi(0)=\psi'(0)=0$, integrating twice yields
$\psi(t)\leq t^2/8$ for positive or negative $t$.
For a fixed system, independence across cases therefore gives
\[
 \acamE\exp\left(t\sum_i(R_m-L_{i,m})\right)
 \leq\exp(nt^2/8),\qquad t>0.
\]
Markov's inequality, optimized at $t=4\gamma$, yields
\[
 \acamP(R_m-\widehat R_m>\gamma)
 \leq\inf_{t>0}e^{-tn\gamma+nt^2/8}
 =e^{-2n\gamma^2}=\delta/M.
\]
A union bound gives simultaneous validity; independence among systems
is unnecessary. Clipping at 1 preserves an upper bound because $R_m\leq1$.
On the simultaneous-validity event, every member of
$\mathcal C_\alpha$ has $R_m\leq U_m\leq\alpha$, regardless of how a
member is selected.
\end{proof}

The guarantee concerns population-average terminal loss, not individual
cases, terminal-pattern subgroups, or selective risk conditional on
answering. Within-case observations and stopping may be dependent:
the independent unit is the complete case. Related edits and repeated
patient images must not be counted as iid cases. If the certified set is
empty, the procedure supplies no certificate. Retraining, unregistered
threshold search, changing prompts, or deployment drift changes the
certified object. For $M=20,\delta=0.05,n=203$, the margin is about
$0.12148$, so even zero empirical errors cannot certify a $0.10$ target
with this particular bound.

\begin{acamprop}[T6: selection and calibration]
\label{prop:acam-selection}
Marginal calibration does not imply calibration after history-based
selection. However, if $C\in[0,1]$,
$p=\acamE[C\mid\mathcal F]$, and the selection indicator $J$ is
$\mathcal F$-measurable, then
\[
 \acamE[C\mid\sigma(p,J)]=p\quad\text{almost surely}.
\]
\end{acamprop}

\begin{proof}
For the counterexample, acquire $S\sim\operatorname{Bernoulli}(1/2)$
and let a different target concept be $C=S$. Assign that concept the
constant score $p=1/2$. Then $\acamE[C\mid p]=1/2=p$, so marginal
calibration is exact. Select this concept only when $S=1$, using only
previously acquired evidence. On the selected subset,
$\acamE[C\mid p=1/2,J=1]=1$, which differs from the reported score.
The example establishes possibility, not the optimality of querying a
redundant concept.

For the positive claim, both $p$ and $J$ are $\mathcal F$-measurable.
The tower property gives
\[
 \acamE[C\mid\sigma(p,J)]
 =\acamE[\acamE[C\mid\mathcal F]\mid\sigma(p,J)]
 =\acamE[p\mid\sigma(p,J)]=p.
\]
Independent selection randomization can be incorporated into
$\mathcal F$ without changing $\acamE[C\mid\mathcal F]$.
\end{proof}

Thus adaptivity need not invalidate a correct full-information
conditional probability. The concern is that marginal calibration is
weaker than conditioning on all information used to select a query.

\begin{acamprop}[T7: conditional-ranking gap]
\label{prop:acam-conditional-gap}
Suppose a paid first semantic query reveals $U$, followed by one action
from a common finite feasible set $\mathcal A$. Fix the response mechanism,
terminal head, loss, and remaining costs. For each action $a$, let $L_a$
be its integrable terminal loss plus cost and let
$q_a(u)=\acamE[L_a\mid U=u]$. The optimal static and $U$-adaptive
second-action risks are
\[
 R_{\rm stat}=\min_{a\in\mathcal A}\acamE q_a(U),\qquad
 R_{\rm dyn}=\acamE\min_{a\in\mathcal A}q_a(U).
\]
Their gap $G=R_{\rm stat}-R_{\rm dyn}$ is nonnegative, and is zero if and
only if some fixed action is conditionally optimal almost surely. With
two actions, writing $d(U)=q_1(U)-q_2(U)$,
\[
 G=\tfrac12\bigl(\acamE|d(U)|-|\acamE d(U)|\bigr). \tag{T7}
\]
If a deployable estimate satisfies
$\sup_a|\widehat q_a(u)-q_a(u)|\leq\epsilon$ almost surely, its
pointwise minimizer has expected risk at most $R_{\rm dyn}+2\epsilon$.
\end{acamprop}

\begin{proof}
Pointwise, $\min_bq_b(U)\leq q_a(U)$ for every fixed $a$; taking
expectations and then minimizing over $a$ gives $G\geq0$. Let $a^*$
minimize the static risk, which exists because $\mathcal A$ is finite.
If $G=0$, the nonnegative random variable
$q_{a^*}(U)-\min_bq_b(U)$ has zero mean and thus vanishes almost surely.
The converse is immediate. For two actions, use
$\min(v,w)=(v+w-|v-w|)/2$ both before and after taking expectations.
Finally, for a true conditional minimizer $a_u^*$ and the learned
minimizer $\widehat a(u)$, pointwise
\[
 q_{\widehat a}(u)
 \leq\widehat q_{\widehat a}(u)+\epsilon
 \leq\widehat q_{a_u^*}(u)+\epsilon
 \leq q_{a_u^*}(u)+2\epsilon.
\]
Taking expectations proves the last assertion.
\end{proof}

This standard decision identity is a diagnostic, not a novel algorithm
guarantee. The CUB OOF-fit/validation-evaluate two-query pilot shows a
local automatic-concept cross-entropy signal, with mixed accuracy, but
does not estimate the population-optimal gap or establish K16, stopping,
JEV-specific, real-cost, or independent-test gains. Changing candidate
menus requires a separately defined static comparator and assumptions.

\begin{acamprop}[T8: softmax-imitation ranking can invert expected utility]
\label{prop:acam-choice-softmax-gap}
At one fixed controller-visible state, let the legal set contain STOP
and 28 unit-cost singleton queries. For binary realized terminal errors,
define per-case utility $u_a=e_a+\lambda\mathbf{1}\{a\ne\mathrm{STOP}\}$ and the
complete-set teacher $q_a=\exp(-u_a/\theta)/\sum_b\exp(-u_b/\theta)$.
Even with $\lambda=0.03$ and $\theta=0.5$, an exact population minimizer
of cross-entropy to this teacher can select an action with strictly
higher expected utility than another feasible action.
\end{acamprop}

\begin{proof}
Name three queries $A,B,C$ and the other 25 dummy queries. Two case types
are indistinguishable at the fixed visible state. Type I has probability
$p=0.52$, and type II has probability $1-p$. In the coordinate order
$(\mathrm{STOP},A,B,C,\mathrm{dummy})$, their binary error vectors are
\[
 (1,0,1,0,1)\quad\text{and}\quad(1,1,0,1,1),
\]
where every dummy has error one. Thus
$\acamE u_A=\acamE u_C=1-p+\lambda=0.51$ and
$\acamE u_B=p+\lambda=0.55$; STOP and dummy expected utilities are one
and $1+\lambda$. For example, these errors are realizable with hidden
class labels one and two for types I and II, a fixed head predicting
class zero at STOP/after dummy queries, class one after $A$ or $C$, and
class two after $B$; the singleton answers may be constant. No type
indicator enters the controller-visible state.

Put $r=\exp(-1/\theta)$ and $c=\exp(-\lambda/\theta)$. The two teacher
normalizers are
\[
 D_I=2c+26cr+r,\qquad D_{II}=c+27cr+r.
\]
Their difference is $c(1-r)>0$. In particular,
\[
 \acamE(q_A-q_B)=c(1-r)
   \left(\frac{p}{D_I}-\frac{1-p}{D_{II}}\right)<0
\]
because $p=0.52<D_I/(D_I+D_{II})\simeq0.541331$. Direct substitution
gives $\acamE q_B\simeq0.112475>\acamE q_A=\acamE q_C\simeq0.105373$;
STOP and every dummy have smaller mean teacher probability. For a
fixed visible state the expected soft-target cross-entropy is minimized
uniquely by $\pi_a=\acamE q_a$, since every mean probability is positive and
the excess cross-entropy is $D_{\rm KL}(\acamE q\Vert\pi)$. Its argmax
therefore selects $B$, although $A$ or $C$ has expected utility lower
by $0.04$.
\end{proof}

This is a counterexample to a Bayes-optimality interpretation of
one-step softmax Choice imitation, not a claim that the fitted CUB policy
fails or that a separate multi-step Bellman target is inconsistent. It
uses the same STOP/singleton menu, Boolean errors, unit group costs,
$\lambda$, and temperature as the current structured Choice pilot; its
two hidden case types are not asserted to describe CUB.

\paragraph{Scope and primary references.}
The concentration background is Hoeffding (1963),
\href{https://doi.org/10.1080/01621459.1963.10500830}
{Probability Inequalities for Sums of Bounded Random Variables}.
The finite-family protocol is a conservative instance of
\href{https://arxiv.org/abs/2110.01052}{Learn then Test}.
\href{https://arxiv.org/html/2606.16667v4}{BCEA, version 4}
already incorporates acquisition into post-acquisition recalibration;
its strict finite-grid procedure and practical threshold scans must not
be conflated.
\href{https://arxiv.org/html/2608.15520v2}{RouteCert, version 2}
already studies complete policy--pattern validation and the noiseless
XOR obstruction. We make no first-of-its-kind calibration or
complementarity claim. The proof-obligation ledger separately records
unverified empirical assumptions, including learned value accuracy,
semantic fidelity, iid case structure, and real execution costs.
